% ch08 -- Landauer: Phasenraum, nicht Semantik
\label{ch:landauer_phasenraum}
% Quelldokumente: A271, A280, 301, 302

\epigraph{Die Physik zählt Gebiete. Sie liest sie nicht.}{A271}

\section{Liouville und das Gebiet}

Der Satz von Liouville: Das Phasenraumvolumen eines abgeschlossenen Systems
ist invariant. Wenn eine Gebietsabbildung $f: \{A, B\} \to Z$ mehrere
Eingangsgebiete auf ein Zielgebiet reduziert, muss das „fehlende" Volumen
irgendwo hingehen. Das Bad übernimmt es. \status{B}

Das ist der physikalische Kern der Landauer-Grenze --- ohne Semantik,
ohne Informationsbegriff.

\section{Drei Gebietsoperationen}

\textbf{Bijektive Abbildung:} $W_{\text{nachher}} = W_{\text{vorher}}$.
Keine Kosten. Transistor kippt von A nach B --- eine bijektive Operation.
Kein Volumenverlust. \status{B}

\textbf{Nicht-bijektive Abbildung:} $W_{\text{nachher}} < W_{\text{vorher}}$.
Mindestkosten: $Q \ge k_BT \ln(W_{\text{vorher}}/W_{\text{nachher}})$.
RESTORE TO ONE --- das klassische Bit-Löschen. \status{B}

\textbf{Dekohärenz:} Das Superpositions-Gebiet $\{|\psi\rangle\}$ wird auf
ein klassisches Gebiet reduziert. Das ist ebenfalls eine nicht-bijektive
Gebietsabbildung. Die Kosten fallen ans Bad --- genau Landauer. \status{B}

\begin{laierspur}
Bijektiv ist wie eine Umstellung der Möbel --- dieselben Gegenstände,
andere Anordnung, nichts verloren. Nicht-bijektiv ist wie Möbel
zusammenschieben --- eine Ecke des Raums bleibt frei, und der Raum
hat dadurch nicht mehr Platz bekommen: Der frei gewordene Teil
muss sich irgendwo sonst wiederfinden.
\end{laierspur}

\section{Der Umkehrtest}

Dok.\,301 formuliert den \textbf{Umkehrtest} als operationales Kriterium:
Eine Operation ist genau dann Landauer-kostenpflichtig, wenn ihre Umkehrung
\emph{physikalisch} nicht durchführbar ist --- ohne Einbeziehung eines Bads.

Das ist schärfer als die übliche Formulierung: Nicht die logische
Irreversibilität (Landauer 1961) entscheidet, sondern die physikalische
Nicht-Bijektivität der Gebietsabbildung. \status{B Dok.\,301}

\section{Elementarzellen und spektrale Antwort}

Wenn man den Phasenraum in Elementarzellen $h^f$ aufteilt, gibt jede
Zelle genau eine orthogonale Antwort. Die spektrale Struktur der
Landauer-Buchhaltung ist damit vollständig durch $h^f$ und das Verhältnis
$W_{\text{vorher}}/W_{\text{nachher}}$ bestimmt --- keine freien Parameter.
\status{B Dok.\,302}
